\documentclass[8pt,a4paper,twocolumn]{article}
\usepackage[utf8]{inputenc}
\usepackage[spanish]{babel}
\usepackage[margin=0.6cm]{geometry}
\usepackage{listings}
\usepackage{xcolor}
\usepackage{enumitem}

% Configuración de colores para código
\definecolor{codegreen}{rgb}{0,0.6,0}
\definecolor{codegray}{rgb}{0.5,0.5,0.5}
\definecolor{codepurple}{rgb}{0.58,0,0.82}
\definecolor{backcolour}{rgb}{0.95,0.95,0.92}

\lstdefinestyle{sqlstyle}{
    backgroundcolor=\color{backcolour},
    commentstyle=\color{codegreen},
    keywordstyle=\color{blue},
    numberstyle=\tiny\color{codegray},
    stringstyle=\color{codepurple},
    basicstyle=\ttfamily\tiny,
    breakatwhitespace=false,
    breaklines=true,
    captionpos=b,
    keepspaces=true,
    showspaces=false,
    showstringspaces=false,
    showtabs=false,
    tabsize=2
}

\lstset{
    basicstyle=\ttfamily\tiny,
    keywordstyle=\color{blue}\bfseries,
    commentstyle=\color{green},
    stringstyle=\color{red},
    showstringspaces=false,
    breaklines=true,
    frame=single,
    numbers=left,
    numberstyle=\tiny\color{gray},
    inputencoding=utf8,
    literate={á}{{\'a}}1 {é}{{\'e}}1 {í}{{\'\i}}1 {ó}{{\'o}}1 {ú}{{\'u}}1
             {Á}{{\'A}}1 {É}{{\'E}}1 {Í}{{\'I}}1 {Ó}{{\'O}}1 {Ú}{{\'U}}1
             {ñ}{{\~n}}1 {Ñ}{{\~N}}1 {ü}{{\"u}}1 {Ü}{{\"U}}1
}

\lstset{style=sqlstyle}

\setlist[itemize]{noitemsep, topsep=0pt, leftmargin=*, itemsep=0pt}
\setlist[enumerate]{noitemsep, topsep=0pt, leftmargin=*, itemsep=0pt}
\setlength{\columnsep}{0.6cm}
\setlength{\parindent}{0pt}
\setlength{\parskip}{0.5pt}

\title{\textbf{Tema 3: PL/SQL Básico}}
\author{}
\date{}

\begin{document}

\tiny

\maketitle

\section*{1. Bloques PL/SQL}

\textbf{Estructura}: DECLARE (opcional), BEGIN (obligatorio), EXCEPTION (opcional).

\begin{lstlisting}
DECLARE
  v_nom VARCHAR2(50); v_sal NUMBER := 0;
BEGIN
  SELECT nom,sal INTO v_nom,v_sal FROM emp WHERE id=101;
  DBMS_OUTPUT.PUT_LINE('Emp: ' || v_nom);
EXCEPTION
  WHEN NO_DATA_FOUND THEN DBMS_OUTPUT.PUT_LINE('No existe');
  WHEN OTHERS THEN DBMS_OUTPUT.PUT_LINE('Error: '||SQLERRM);
END;
/
\end{lstlisting}

\textbf{Tipos}: Anónimo (sin nombre), Nombrado (proc/func/trigger), Anidado.

\textbf{Anidado}:
\begin{lstlisting}
DECLARE v_ext NUMBER := 10;
BEGIN
  DECLARE v_int NUMBER := 20;
  BEGIN
    NULL; -- v_ext y v_int disponibles
  END;
  -- v_int no existe aquí
END;
/
\end{lstlisting}

\section*{2. Variables}

\textbf{Tipos}: NUMBER, VARCHAR2(n), CHAR(n), DATE, BOOLEAN, TIMESTAMP.

\begin{lstlisting}
DECLARE
  v_nom VARCHAR2(100); v_edad NUMBER(3);
  v_sal NUMBER(10,2) := 3000.50; v_activo BOOLEAN := TRUE;
  v_fecha DATE := SYSDATE; c_iva CONSTANT NUMBER := 0.21;
BEGIN NULL; END;
/
\end{lstlisting}

\textbf{\%TYPE} (tipo de columna):
\begin{lstlisting}
v_nom emp.nom%TYPE; v_sal emp.sal%TYPE;
\end{lstlisting}

\textbf{\%ROWTYPE} (fila completa):
\begin{lstlisting}
DECLARE
  v_emp emp%ROWTYPE;
BEGIN
  SELECT * INTO v_emp FROM emp WHERE id=101;
  DBMS_OUTPUT.PUT_LINE(v_emp.nom||': '||v_emp.sal);
END;
/
\end{lstlisting}

\textbf{Asignación}: \texttt{v\_nom := 'Juan'; v\_i := v\_i + 1;}

\section*{3. Control}

\textbf{IF-THEN-ELSE}:
\begin{lstlisting}
IF v_sal>5000 THEN v_cat:='Alto';
ELSIF v_sal>3000 THEN v_cat:='Medio';
ELSE v_cat:='Bajo'; END IF;
\end{lstlisting}

\textbf{CASE}:
\begin{lstlisting}
CASE v_dia
  WHEN 1 THEN DBMS_OUTPUT.PUT_LINE('Lunes');
  WHEN 2 THEN DBMS_OUTPUT.PUT_LINE('Martes');
  ELSE DBMS_OUTPUT.PUT_LINE('Otro'); END CASE;
\end{lstlisting}

\textbf{LOOP}:
\begin{lstlisting}
LOOP
  v_i := v_i + 1;
  EXIT WHEN v_i > 10;
END LOOP;
\end{lstlisting}

\textbf{WHILE/FOR}:
\begin{lstlisting}
WHILE v_i <= 5 LOOP v_i:=v_i+1; END LOOP;
FOR i IN 1..10 LOOP DBMS_OUTPUT.PUT_LINE(i); END LOOP;
FOR i IN REVERSE 1..5 LOOP NULL; END LOOP;
\end{lstlisting}

\section*{4. Cursores}

\textbf{Implícito (1 fila)}:
\begin{lstlisting}
SELECT nom,sal INTO v_nom,v_sal FROM emp WHERE id=101;
DBMS_OUTPUT.PUT_LINE('Filas: ' || SQL%ROWCOUNT);
\end{lstlisting}

\textbf{Explícito (múltiples)}:
\begin{lstlisting}
DECLARE
  CURSOR c IS SELECT nom,sal FROM emp WHERE dept=10;
  v_nom emp.nom%TYPE; v_sal emp.sal%TYPE;
BEGIN
  OPEN c;
  LOOP
    FETCH c INTO v_nom,v_sal;
    EXIT WHEN c%NOTFOUND;
    DBMS_OUTPUT.PUT_LINE(v_nom||': '||v_sal);
  END LOOP;
  CLOSE c;
END;
/
\end{lstlisting}

\textbf{FOR LOOP (automático)}:
\begin{lstlisting}
FOR r IN (SELECT nom,sal FROM emp) LOOP
  DBMS_OUTPUT.PUT_LINE(r.nom);
END LOOP;
\end{lstlisting}

\textbf{Atributos}: \%FOUND, \%NOTFOUND, \%ROWCOUNT, \%ISOPEN.

\section*{5. Registros}

\textbf{TYPE RECORD}: Estructura de datos compuesta.

\begin{lstlisting}
DECLARE
  TYPE r_persona IS RECORD(
    nombre VARCHAR2(25),
    apellido VARCHAR2(25),
    salario NUMBER(10));
  v_emp r_persona;
BEGIN
  v_emp.nombre := 'Ana';
  v_emp.salario := 30000;
END;
/
\end{lstlisting}

\textbf{\%ROWTYPE para registros}:
\begin{lstlisting}
DECLARE
  v_emp empleados%ROWTYPE;  -- Estructura de tabla
  CURSOR c IS SELECT nom,sal FROM emp;
  v_dato c%ROWTYPE;          -- Estructura de cursor
BEGIN
  SELECT * INTO v_emp FROM empleados WHERE id=101;
  DBMS_OUTPUT.PUT_LINE(v_emp.nom || ': ' || v_emp.sal);
END;
/
\end{lstlisting}

\textbf{Cursores con registros}:
\begin{lstlisting}
DECLARE
  CURSOR c IS SELECT nom,sal FROM emp;
BEGIN
  FOR v_reg IN c LOOP
    DBMS_OUTPUT.PUT_LINE(v_reg.nom || ': ' || v_reg.sal);
  END LOOP;
END;
/
\end{lstlisting}

\section*{6. Excepciones}

\textbf{Predefinidas}: NO\_DATA\_FOUND, TOO\_MANY\_ROWS, ZERO\_DIVIDE, VALUE\_ERROR, INVALID\_NUMBER.

\begin{lstlisting}
BEGIN
  SELECT sal INTO v_sal FROM emp WHERE id=999;
EXCEPTION
  WHEN NO_DATA_FOUND THEN DBMS_OUTPUT.PUT_LINE('No existe');
  WHEN TOO_MANY_ROWS THEN DBMS_OUTPUT.PUT_LINE('Múltiples');
  WHEN OTHERS THEN DBMS_OUTPUT.PUT_LINE('Error: '||SQLERRM);
END;
/
\end{lstlisting}

\textbf{Usuario}:
\begin{lstlisting}
DECLARE
  edad_inv EXCEPTION;
BEGIN
  IF v_edad<18 THEN RAISE edad_inv; END IF;
EXCEPTION
  WHEN edad_inv THEN DBMS_OUTPUT.PUT_LINE('Edad >=18');
END;
/
\end{lstlisting}

\section*{7. Operaciones SQL}

\textbf{SELECT INTO}:
\begin{lstlisting}
SELECT nombre,salario INTO v_nom,v_sal FROM emp WHERE id=101;
\end{lstlisting}

\textbf{INSERT/UPDATE/DELETE}:
\begin{lstlisting}
INSERT INTO emp VALUES(999,'Ana',3500); COMMIT;
UPDATE emp SET salario=salario*1.1 WHERE dept=10;
DELETE FROM emp WHERE salario<1000;
DBMS_OUTPUT.PUT_LINE('Filas: ' || SQL%ROWCOUNT);
\end{lstlisting}

\section*{8. DBMS\_OUTPUT}

\begin{lstlisting}
SET SERVEROUTPUT ON;
BEGIN
  DBMS_OUTPUT.PUT_LINE('Mensaje');
END;
/
\end{lstlisting}

\section*{9. Ejemplos Prácticos}

\textbf{Calcular bonus}:
\begin{lstlisting}
DECLARE
  CURSOR c IS SELECT emp_id,salario FROM empleados WHERE dept=10;
BEGIN
  FOR r IN c LOOP
    UPDATE empleados SET bonus=r.salario*0.1 WHERE emp_id=r.emp_id;
  END LOOP;
  COMMIT;
END;
/
\end{lstlisting}

\textbf{Salario mínimo}:
\begin{lstlisting}
DECLARE
  c_min CONSTANT NUMBER := 1500;
BEGIN
  UPDATE empleados SET salario=c_min WHERE salario<c_min;
  DBMS_OUTPUT.PUT_LINE('Actualizados: ' || SQL%ROWCOUNT);
  COMMIT;
END;
/
\end{lstlisting}

\end{document}
